%%
%% fall-lecture04.tex
%% 
%% Made by Alex Nelson <pqnelson@gmail.com>
%% Login   <alex@lisp>
%% 
%% Started on  2026-10-08T07:54:02-0700
%% Last update 2026-10-08T07:54:02-0700
%% 

\lecture{}

\begin{node}[Conservation of linear momentum]
Taking $b=\vec{u}$ a vector (or, the components of the vector),
Newton's second Law tells us that the force on the system $\vec{F}$
satisfies
\begin{equation}
\vec{F}=\frac{\D\vec{p}}{\D t}=\frac{D\vec{p}}{D t}=\frac{D}{Dt}\iiint_{V_{s}}\rho\vec{u}\,\D V_{s}
\end{equation}
Now we can use Reynolds transport theorem to get
\begin{equation}
\vec{F}=\frac{\partial}{\partial t}\iiint_{V_{c}}\rho\vec{u}\,\D V_{c}
-\oint_{\boundary V_{c}}\rho\vec{u}(\vec{u}\cdot\UnitNormalVec{n})\,\D A_{c}
\end{equation}
where we treat $\vec{F}$ as the force acting on the control volume,
since we are stipulating the initial conditions that
$V_{c}(t_{0})=V_{s}(t_{0})$ initially.
\end{node}

\begin{definition}
We decompose force $\vec{F}$ into a sum of \define{Surface Forces}
(e.g., pressure forces, viscuous forces, etc.) and \define{Body Forces}
(e.g., gravity, etc.)
\begin{equation}
\vec{F}=\vec{F}_{\text{surface}}+\vec{F}_{\text{body}}.
\end{equation}
The surface forces are written as
\begin{equation}
\vec{F}_{\text{surface}}=\oint_{\boundary V_{c}}\vec{f}\,\D A_{c}
\end{equation}
where $\vec{f}$ is the force per unit area, and the body forces are
written as
\begin{equation}
\vec{F}_{\text{body}}=\iiint_{V_{c}}\vec{\gamma}\,\D V_{c}.
\end{equation}
For example, $\vec{\gamma}=\rho\vec{g}$ for gravitational forces.
\end{definition}

\begin{example}
Suppose we have a chamber with an inlet and an outlet (a
``nozzle''). Our control volume is drawn in dashed red line:
\begin{center}
\includegraphics{img/img.2}
\end{center}
We assume:
\begin{enumerate}
\item incompressibility ($\rho$ is constant),
\item we can neglect wall stress (from ``hitting'' the wall),
\item we can ignore gravity.
\end{enumerate}
We have mass conservation gives us
\begin{subequations}
  \begin{align}
0 &= \oint_{\boundary V_{c}}\rho\vec{u}\cdot\UnitNormalVec{n}\,\D A_{c}\\
&=\oint_{\text{inlet}}\rho\vec{u}\cdot\UnitNormalVec{n}\,\D A_{c} + \oint_{\text{wall}}\rho\vec{u}\cdot\UnitNormalVec{n}\,\D A_{c} + \oint_{\text{outlet}}\rho\vec{u}\cdot\UnitNormalVec{n}\,\D A_{c}\\
&=\oint_{\text{inlet}}\rho\vec{u}\cdot\UnitNormalVec{n}\,\D A_{c} + 0 + \oint_{\text{outlet}}\rho\vec{u}\cdot\UnitNormalVec{n}\,\D A_{c}\\
&=-\rho_{1}A_{1}u_{1}+\rho_{2}A_{2}u_{2}\\
&=\rho(-A_{1}u_{1}+A_{2}u_{2}),
  \end{align}
\end{subequations}
which implies
\begin{equation}
A_{1}u_{1}=A_{2}u_{2}.
\end{equation}
Now, we can use the conservation of linear momentum to write
\begin{subequations}
  \begin{align}
\vec{F} &= \mbox{(inlet)} + \begin{pmatrix}\mbox{pressure forces from}\\
  \mbox{wall adjacent to}\\
  \mbox{inlet}
\end{pmatrix} + \begin{pmatrix}\mbox{forces from}\\
  \mbox{outlet}
\end{pmatrix}\\
&= -p_{1}A_{1}-p_{1}(A_{2}-A_{1})+p_{2}A_{2}\\
&= -\rho u_{1}^{2}A_{1} + \rho u_{2}^{2}A_{2}\quad\mbox{by RTT},
\end{align}
\end{subequations}
which gives us
\begin{equation}
p_{2}-p_{1}=\rho u_{2}(u_{2}-u_{1}).
\end{equation}
Although this is a highly simplified and highly idealized situation,
we can remove each assumption to obtain successive approximations
until we get a suitable approximation.
\end{example}

\begin{node}[Conservation of angular momentum]
Schematically, we have the following situation:
\begin{center}
\includegraphics{img/img.3}
\end{center}
We have an origin $O$, we have a control volume CV. Then we set
$b=\vec{r}\times\vec{u}$ to track the angular momentum. We use Euler's
equations of motion (the generalization of Newton's equations for
angular momentum)
\begin{equation}
\vec{M} = \frac{D}{D t}\iiint_{V_{s}}\rho\vec{r}\times\vec{u}\,\D V_{s}
\end{equation}
where $\vec{M}$ is the torque. Then
\begin{equation}
\vec{M}=\frac{\partial}{\partial t}\iiint_{V_{c}}\rho\vec{r}\times\vec{u}\,\D V_{c}
+ \oint_{\boundary V_{c}}\rho(\vec{r}\times\vec{u})(\vec{u}\cdot\UnitNormalVec{n})\,\D A_{c}.
\end{equation}
We can partition the torque $\vec{M}$ into the sum of surface torques
and body torques
\begin{equation}
\vec{M}=\vec{M}_{\text{surface}}+\vec{M}_{\text{body}}
\end{equation}
where
\begin{equation}
\vec{M}_{\text{surface}}=\oint_{\boundary V_{c}}\vec{r}\times\vec{f}\,\D A_{c}
\end{equation}
and
\begin{equation}
\vec{M}_{\text{body}}=\iiint_{V_{c}}\vec{r}\times\vec{\gamma}\,\D V_{c}
\end{equation}
using the surface and body force densities.
\end{node}

\begin{node}[Conservation of energy]
We have
\begin{subequations}
  \begin{align}
b &=
\begin{pmatrix}\mbox{kinetic}\\
\mbox{energy}
\end{pmatrix} + \begin{pmatrix}
\mbox{specific internal}\\
\mbox{energy}
\end{pmatrix}\\
&=\frac{1}{2}\vec{u}\cdot\vec{u} + e\\
&=\varepsilon
  \end{align}
\end{subequations}
where $\varepsilon$ is the energy per unit mass. The we see the volume
integral over the system
\begin{equation}
B = \iiint_{V_{s}}\rho b\,\D V_{s} = E.
\end{equation}
Then we use Thermodynamics to write
\begin{subequations}
  \begin{align}
\begin{pmatrix}\mbox{change in energy}\\
\mbox{stored in the system}
\end{pmatrix} &= \begin{pmatrix}\mbox{Heat transfer}\\
\mbox{to system by}\\
\mbox{environment}
\end{pmatrix} + \begin{pmatrix}\mbox{work done by}\\
  \mbox{surroundings to}\\
  \mbox{system}
\end{pmatrix}\\
\D E &= \D Q + \D W.
\end{align}
\end{subequations}
Then we translate to fluid mechanics using the Reynolds transport
theorem
\begin{subequations}
  \begin{align}
\frac{D E}{D t} &= \frac{D Q}{D t} + \frac{D W}{D t}\\
&=\frac{\partial}{\partial t}\iiint_{V_{c}}\rho\varepsilon\,\D V_{c}
+\oint_{\boundary V_{c}}\rho\varepsilon(\vec{u}\cdot\UnitNormalVec{n})\,\D A_{c}.
  \end{align}
\end{subequations}
We can say a little about the work done, since it is related to force:
\begin{equation}
\frac{D W}{D t} = \iiint_{V_{c}}\vec{\gamma}\cdot\vec{u}\,\D V_{c}
+\oint_{\boundary V_{c}}\vec{f}\cdot\vec{u}\,\D A_{c}.
\end{equation}
Similarly, the heat transfer can be written as the sum
\begin{equation}
\frac{D Q}{D t} = \iiint_{V_{c}}\rho \mathcal{Q}\,\D V_{c}
-\oint_{\boundary V_{c}}\vec{q}\cdot\UnitNormalVec{n}\,\D A_{c},
\end{equation}
where $\vec{q}=-\kappa\vec{\nabla}T$ is the heat flux vector, and
$\mathcal{Q}$ is the heat source strength per unit mass.
The heat contributions are exogeneous quantities given as a sort of
``source terms'' to the equations.

Altogether, we get
\begin{equation}
\iiint\vec{\gamma}\cdot\vec{u}+\rho\mathcal{Q}\,\D V_{c}
+\oint_{\boundary V_{c}}\vec{f}\cdot\vec{u}-\vec{q}\cdot\UnitNormalVec{n}\,\D A_{c}
=\frac{\partial}{\partial t}\iiint_{V_{c}}\rho\varepsilon\,\D V_{c}
+\oint_{\boundary V_{c}}\rho\varepsilon(\vec{u}\cdot\UnitNormalVec{n})\,\D A_{c}.
\end{equation}
This is the conservation of energy.
\end{node}

\begin{node}[Recap]
We used the Reynolds transport theorem \emph{and} some physics to
obtain the integral conservation equations of mass, momnetum, and
energy for finite control volumes. We looked at 4 different
properties: mass, linear momentum, angular momentum, and energy. Next
time, we will transform these into differential form, a first step
towards the Navier--Stokes equations.
\end{node}